% 2.3 章 · 正交多项式与函数逼近
% 这一章的笔记要把"用多项式逼近一个函数"拆成四件互不相同的数学事实：
%   1) 正交多项式怎么造、零点在哪里（三项递推 + Sturm 计数）；
%   2) 节点和权重怎么选才"免费"拿到 2n-1 阶精度（Gauss 求积）；
%   3) "平方意义下最好"是什么（正交投影，正规方程）；
%   4) "最大误差最小"是什么（Chebyshev 交错定理 + de la Vallee Poussin 下界）；
% 再加一节"误差衰减得有多快"：光滑性 -> 收敛率，解析函数走 Bernstein 椭圆。
% 文中每一个数字都是 tests/2.3 与 server/chapters/2.3.py 里跑出来的实测值。
% 断言附表由 tools/build_book.py 从实际运行结果生成，\input 进来，不手抄。

\numlabchapter{正交多项式与函数逼近}{2.3}

\section{这一章要解决什么}

给定 $f\in C[-1,1]$ 与次数 $n$，我们要找"最好的"多项式。麻烦在于"最好"至少有三种意思：

\begin{enumerate}
  \item \textbf{在节点上相等}：插值。它不告诉你误差有多大，节点选坏了还会发散。
  \item \textbf{平方误差最小}：$L^2$ 投影。它把问题变成一次对称正定的线性方程组
        （正规方程），并且残差与整个子空间正交——这是"投影"这个词的几何含义。
  \item \textbf{最大误差最小}：$L^\infty$（一致）最佳逼近。它的判别法是
        \textbf{交错定理}：误差曲线必须在 $n+2$ 个点上等幅、符号交替地振荡。
\end{enumerate}

三种"最好"给出三个不同的多项式，而工程上最常被混为一谈的正是后两种。
本章的四个可视化分别对应：投影的几何（\texttt{projection}）、三种光滑性的收敛率
（\texttt{convergence\_rates}）、交错点的样子（\texttt{equioscillation}）、
以及"节点从哪来"（\texttt{zeros\_evolution}）。

记号：$u=\tfrac12\varepsilon_{\text{mach}}=2^{-53}\approx1.11\times10^{-16}$ 是单位舍入；
$\langle f,g\rangle_w=\int_I f g\,w$ 是带权内积，$\|f\|_w^2=\langle f,f\rangle_w$；
$\Phi_n=\mathrm{span}\{p_0,\dots,p_n\}$；$\|\cdot\|_\infty$ 是区间上的最大模。

\section{正交多项式：三项递推与零点}

\begin{definition}[正交多项式]
设 $I\subseteq\mathbb{R}$，权函数 $w>0$ 在 $I$ 上可积。若 $\deg p_k=k$ 且
\[
  \langle p_i,p_j\rangle_w=0\qquad(i\ne j),
\]
则称 $\{p_k\}$ 是关于 $w$ 的正交多项式族。若还有 $\|p_k\|_w=1$，称为\textbf{正交归一}族，
记作 $\hat p_k$。
\end{definition}

\begin{theorem}[三项递推]
\label{thm:recurrence}
设 $\{p_k\}$ 是首一正交多项式族，$h_k=\|p_k\|_w^2$。则存在常数 $\alpha_k,\beta_k$ 使
\[
  p_{k+1}(x)=(x-\alpha_k)p_k(x)-\beta_k\,p_{k-1}(x),\qquad
  \alpha_k=\frac{\langle xp_k,p_k\rangle_w}{h_k},\quad
  \beta_k=\frac{h_k}{h_{k-1}},
\]
且 $\beta_k>0$（$k\ge1$）。反之，任何满足该递推、$\beta_k>0$ 的多项式列都正交。
\end{theorem}

\begin{proof}
$x p_k$ 是 $k+1$ 次多项式，把它按正交基 $\{p_0,\dots,p_{k+1}\}$ 展开：
\[
  x p_k=\sum_{j=0}^{k+1} c_j p_j,\qquad
  c_j=\frac{\langle xp_k,p_j\rangle_w}{h_j}.
\]
对 $j\le k-2$，$\langle xp_k,p_j\rangle_w=\langle p_k,xp_j\rangle_w=0$，因为 $xp_j$ 的次数
$\le k-1$，与 $p_k$ 正交。于是只剩 $c_k,c_{k-1},c_{k+1}$ 三项。比较首项系数：
$xp_k$ 与 $p_{k+1}$ 都是首一 $k+1$ 次多项式，故 $c_{k+1}=1$，即
\[
  p_{k+1}=xp_k-c_kp_k-c_{k-1}p_{k-1},
\]
这正是要证的递推，其中 $c_k=\alpha_k$、$c_{k-1}=\beta_k$。由 $\langle xp_k,p_{k-1}\rangle_w
=\langle p_k,xp_{k-1}\rangle_w=h_k$（$xp_{k-1}=p_k+\text{低次}$）得 $\beta_k=h_k/h_{k-1}>0$。
反向命题：若递推成立，对 $k>j$ 用归纳法算 $\langle p_{k+1},p_j\rangle_w$，
低次项被正交性消掉，只剩下 $\beta_k\langle p_{k-1},p_j\rangle_w$ 与
$\alpha_k$ 项，逐层归零。
\end{proof}

\textbf{实测}（\texttt{test\_recurrence\_matches\_analytic\_forms}）：自研递推与闭式解的最大
绝对差，Legendre 在 $n\le12$ 上是 $3.0\times10^{-13}$，Chebyshev 是 $3.7\times10^{-15}$。
Laguerre 必须限定窗口：$x\in[0,5]$ 上 $6.6\times10^{-12}$，$x\in[0,20]$ 上涨到
$1.3\times10^{-7}$——那不是实现不准，而是 $|L_{40}(142)|\sim10^{14}$，
多项式自身量级被放大的结果。这条"绝对误差要看尺度"的教训会在零点残差上再出现一次。

\begin{corollary}[零点的位置与交错]
\label{cor:zeros}
$p_n$ 的 $n$ 个零点全为单实根、全部落在 $I$ 的内部；并且 $p_n$ 与 $p_{n-1}$ 的零点
严格交错：$p_{n-1}$ 的每个零点区间内恰有 $p_n$ 的一个零点。
\end{corollary}

\begin{proof}
设 $x_1<\dots<x_m$ 是 $p_n$ 在 $I^\circ$ 内的全部变号零点，令
$q(x)=\prod_{i=1}^m(x-x_i)$。若 $m<n$，则 $p_nq$ 在 $I$ 上不变号（每个变号零点都被抵消），
于是 $\langle p_n,q\rangle_w\ne0$；但 $q$ 的次数 $<n$，与 $p_n$ 正交，矛盾。故 $m\ge n$。
又 $\deg p_n=n$ 至多有 $n$ 个根，于是 $m=n$ 且没有复根、没有重根（重根不变号，
会被上面的论证漏掉，而 $m=n$ 已经用完了全部根）。交错性：把 $\langle p_n,p_{n-1}\rangle_w=0$
写成 $\int_I p_np_{n-1}w=0$；若某个区间 $(x_i,x_{i+1})$ 内没有 $p_{n-1}$ 的零点，
则 $p_np_{n-1}w$ 在该区间不变号，逐段求和可知积分不可能为零。
\end{proof}

\textbf{实测}（\texttt{test\_roots\_of\_successive\_degrees\_interlace}）：三族、$n\le24$、
共 240 个"相邻低阶零点区间"里，高阶零点个数与 1 的偏离恒为 $0$。

\begin{theorem}[Sturm 变号计数]
\label{thm:sturm}
设 $\{p_k\}$ 满足定理 \ref{thm:recurrence} 的递推且 $\beta_k>0$。令 $V(x)$ 为序列
$p_0(x),p_1(x),\dots,p_n(x)$（去掉零元后）的符号变号次数。则 $V$ 是阶梯函数，且
\[
  V(x)=\#\{x_k:\ x_k \text{ 落在 } x \text{ 与 } b \text{ 之间}\},
\]
其中 $b$ 是"所有 $p_k$ 同号"的那个端点（Legendre/Chebyshev 在 $[-1,1]$ 上取 $b=1$，
Laguerre 在 $[0,\infty)$ 上取 $b=0$）。
\end{theorem}

\begin{proof}
对 $k$ 归纳。$k=1$ 时 $p_1=x-\alpha_0$ 只有一个零点 $\alpha_0$，而 $V$ 只在 $x=\alpha_0$
处变号一次，两者一致。设结论对 $k$ 成立，考察 $p_{k+1}$。由递推式，
$p_{k+1}(x)=0$ 时 $p_{k+1}$ 与 $p_k$ 之间没有变号（都把 $p_{k-1}$ 项顶掉）：
在 $p_{k+1}$ 的零点处 $p_k\ne0$（否则由递推 $p_{k-1}$ 也 $=0$，与相邻两项无公共零点矛盾）。
又 $p_{k+1}$ 的相邻零点之间必有一个 $p_k$ 的零点（对 $\langle p_{k+1},p_k\rangle_w=0$
用本节推论的交错论证），于是每次穿过 $p_{k+1}$ 的零点，$V$ 恰好增加 $1$，
与 $p_{k+1}$ 的零点数同步。
\end{proof}

为什么这条定理重要：它把"求第 $k$ 个根"变成"二分一个整数计数器"。
\textbf{本模块就是靠它把 $n$ 个根隔离到互不相交的区间里}，再用 Newton 抛光——
实测 Laguerre 的经典渐近初值在 $n=3$ 时就会让两个初值落进同一个根的吸引域
（权重和从 $1$ 变成 $2.68$），而隔离之后$\sum w_k$ 的误差是 $10^{-15}$ 量级。
方向必须按族选对：实测 Legendre 上 $V(x)=\#\{x_k>x\}$ 在 33 个探针点上逐点相等，
Laguerre 上则是 $V(x)=\#\{x_k<x\}$；把方向搞反在"零点没越过探针点"时完全不报错。

\section{Gauss 求积：$2n-1$ 阶精度是免费的}

\begin{theorem}[Gauss 求积]
\label{thm:gauss}
设 $x_1,\dots,x_n$ 是 $p_n$ 的零点，令
\[
  w_k=\int_I \ell_k(x)\,w(x)\,dx,\qquad
  \ell_k(x)=\prod_{j\ne k}\frac{x-x_j}{x_k-x_j}.
\]
则 $w_k>0$，且对任意次数 $\le 2n-1$ 的多项式 $q$ 有
\[
  \int_I q(x)w(x)\,dx=\sum_{k=1}^n w_k\,q(x_k).
\]
\end{theorem}

\begin{proof}
\textbf{正性}：$\ell_k^2$ 的次数是 $2n-2\le2n-1$（先假定求积公式对 $\le 2n-1$ 成立，
这在下面独立证明），$\ell_k^2\ge0$ 且 $\ell_k(x_k)=1$，而 $\ell_k^2$ 在节点处只有
$k$ 点非零，于是
\[
  w_k=\int \ell_k^2 w=\sum_{j}w_j\ell_k^2(x_j)=w_k\cdot1>0 .
\]
（也可直接由 $\ell_k^2$ 的插值恒等式得到 $w_k=\int\ell_k^2w>0$。）

\textbf{精度}：任一 $q$ 的次数 $\le2n-1$，做带余除法 $q=p_n\,s+r$，其中
$\deg s\le n-1$、$\deg r\le n-1$。于是
\[
  \int q w=\int p_n s\,w+\int r w=0+\int r w,
\]
因为 $s\in\Phi_{n-1}$ 与 $p_n$ 正交。右端 $\sum_k w_kq(x_k)=\sum_k w_k r(x_k)$
（$p_n(x_k)=0$ 把 $s$ 项杀掉）。而 $r$ 的次数 $\le n-1$，$n$ 个节点的插值公式
（Lagrange 插值是精确的）给出 $\int r w=\sum_k w_k r(x_k)$。两个方向合起来即得结论。
\end{proof}

\begin{theorem}[$2n-1$ 是硬边界]
\label{thm:sharp}
对 $q=p_n^2$（次数 $2n$）有
\[
  \sum_{k=1}^n w_k\,p_n^2(x_k)=0\ne\int_I p_n^2(x)\,w(x)\,dx=h_n>0 .
\]
因此节点数 $n$ 的求积公式不可能对所有 $2n$ 次多项式精确。
\end{theorem}

\begin{proof}
$p_n(x_k)=0$ 使左端为 $0$，而右端是 $\|p_n\|_w^2=h_n>0$。
\end{proof}

\textbf{实测}（\texttt{test\_gauss\_rules\_integrate\_polynomials\_to\_the\_exact\_degree}）：
$k\le2n-1$ 时三族的矩误差最大 $8.9\times10^{-16}$（Laguerre 用相对误差，$3.2\times10^{-14}$）；
$k=2n$ 处的相对误差立刻变成 $1.4\times10^{-2}$（Laguerre $n=4$）。
这条"既要精确、又要在 $2n$ 处当场失败"的双向断言，比"误差很小"有力得多：
它同时排除了"权重公式抄错一个因子"和"节点数写错"。

常用的权重闭式（可由 Christoffel--Darboux 恒等式推出）：
\[
  w_k^{\text{Leg}}=\frac{2}{(1-x_k^2)\,[P_n'(x_k)]^2},\qquad
  w_k^{\text{Lag}}=\frac{x_k}{n^2\,[L_{n-1}(x_k)]^2},\qquad
  w_k^{\text{Cheb}}=\frac{\pi}{n}.
\]
\textbf{实测}：用求根给出的权重与上面两个闭式的最大相对差是 $2.0\times10^{-16}$
（Legendre）与 $2.8\times10^{-16}$（Laguerre）——两条独立路径互相印证。
注意 Chebyshev 的权重与节点无关（常数 $\pi/n$），这是它被当作"手算基准"的原因。

\section{投影：平方意义下的最好}

\begin{theorem}[正规方程与投影]
\label{thm:projection}
设 $\{\varphi_0,\dots,\varphi_{m-1}\}$ 线性无关，$G_{ij}=\langle\varphi_i,\varphi_j\rangle_w$，
$b_i=\langle f,\varphi_i\rangle_w$。则
\[
  P f=\sum_{i} c_i\varphi_i,\qquad Gc=b
\]
是 $f$ 在 $\Phi=\mathrm{span}\{\varphi_i\}$ 上的正交投影，且
\[
  \|f-Pf\|_w=\min_{q\in\Phi}\|f-q\|_w .
\]
此外 $G$ 对称正定，$P^2=P$，残差 $f-Pf\perp\Phi$。
\end{theorem}

\begin{proof}
\textbf{正定}：对 $c\ne0$，$c^{\mathsf T}Gc=\|\sum_ic_i\varphi_i\|_w^2>0$，因为基线性无关。

\textbf{法方程}：设 $q=\sum q_i\varphi_i$。展开
\[
  \|f-q\|_w^2=\|f\|_w^2-2\,q^{\mathsf T}b+q^{\mathsf T}Gq .
\]
这是 $q$ 的严格凸二次函数，驻点唯一，满足 $Gq=b$，即 $q=Pf$。

\textbf{正交与幂等}：$q=Pf$ 时对每个 $j$，
$\langle f-Pf,\varphi_j\rangle_w=b_j-(Gc)_j=0$，即残差与 $\Phi$ 正交。
于是对任意 $q'\in\Phi$，$\langle f-Pf,q'\rangle_w=0$，由勾股定理
$\|f-q'\|_w^2=\|f-Pf\|_w^2+\|Pf-q'\|_w^2\ge\|f-Pf\|_w^2$，取 $q'=Pf$ 取等。
最后，$Pf\in\Phi$ 且 $P$ 在 $\Phi$ 上是恒等映射，故 $P^2=P$。
\end{proof}

在正交归一基下 $G=I$，$c_i=\langle f,\hat p_i\rangle_w$，正规方程连解都不用解。
换到非正交基（例如单项式基 $1,x,x^2,\dots$），$G$ 变成 Hilbert 型矩阵，
$\kappa(G)$ 指数增长：\textbf{实测}单项式基 $n=12$ 时 $\kappa\approx1.0\times10^8$，
$n=24$ 时 $\kappa\approx1.95\times10^{17}$，而 Legendre 基始终是 $1.000$。
后果不是"精度差一点"：$n=24$ 时单项式基算出的投影函数与良态基的差是 $1.2\times10^{-3}$，
而\textbf{正规方程的残差只有 $1.0\times10^{-11}$}。
$\kappa\,u\approx21.7>1$——系数已经全是噪声，残差却依然"很小"。

\section{一致逼近：交错定理与它的证书}

\begin{theorem}[Chebyshev 交错定理]
\label{thm:alternation}
设 $f\in C[a,b]$，$n\ge0$。$p\in\Phi_n$ 是 $f$ 的最佳一致逼近，当且仅当存在
$n+2$ 个点 $a\le x_0<x_1<\dots<x_{n+1}\le b$ 使
\[
  f(x_j)-p(x_j)=\sigma\,(-1)^j\,E,\qquad E=\|f-p\|_\infty>0,\quad \sigma\in\{\pm1\}.
\]
\end{theorem}

\begin{proof}
\textbf{必要性}：设 $p$ 最佳，$E>0$。若不存在这样的 $n+2$ 点集，则存在一个"改进方向"。
把误差曲线 $e=f-p$ 的极值点集分成正负两组；由假设，任何 $n+2$ 个交替点组都不等幅，
于是可以构造一个非零多项式符号一致的扰动：具体地，取 $e$ 的极大值点 $M$
（$e>0$ 的极值）与极小值点 $m$（$e<0$ 的极值），若 $|e|$ 在它们上的最大值集合
含有少于 $n+2$ 个交替点，则可取 $\tilde p\in\Phi_n$ 使得 $\tilde p$ 在 $e>0$
的极值点处 $>0$、在 $e<0$ 的极值点处 $<0$（这是 $n+1$ 个符号约束下的多项式插值问题，
$n+2$ 个交替点正是它无解/有解的临界条件）。于是 $f-(p+\epsilon\tilde p)$ 在这些极值点上的
最大模对充分小的 $\epsilon>0$ 严格小于 $E$，与最优性矛盾。

\textbf{充分性}：设存在这样的 $n+2$ 点，并设 $q\in\Phi_n$ 更好，即
$\|f-q\|_\infty<E$。令 $r=q-p\in\Phi_n$。在第 $j$ 个交错点上，
$r(x_j)=q(x_j)-p(x_j)=(f(x_j)-p(x_j))-(f(x_j)-q(x_j))$，其符号由
$\sigma(-1)^jE$ 决定（因为 $|f-q|<E$ 不能翻转它），于是 $r$ 在
$x_0,\dots,x_{n+1}$ 上符号交替 $n+2$ 次。连续函数符号交替 $n+2$ 次至少要有
$n+1$ 个零点，而 $r$ 是次数 $\le n$ 的非零多项式（若 $r\equiv0$ 则 $q=p$ 并不更好），
矛盾。
\end{proof}

\begin{corollary}[de la Vall\'ee Poussin 下界]
\label{cor:vdp}
若 $q\in\Phi_n$ 与点集 $x_0<\dots<x_{n+1}$ 上的误差符号交替，则
\[
  \min_j|f(x_j)-q(x_j)|\ \le\ \|f-p^*\|_\infty,
\]
其中 $p^*$ 是 $f$ 的最佳一致逼近。
\end{corollary}

\begin{proof}
记 $h=\min_j|f(x_j)-q(x_j)|$。若 $\|f-p^*\|_\infty<h$，则 $r=p^*-q\in\Phi_n$ 在上述
$n+2$ 个点上符号仍然交替（$|f-p^*|\le\|f-p^*\|_\infty<h\le|f-q|$ 不足以翻转符号），
于是 $r$ 至少有 $n+1$ 个零点，只能 $r\equiv0$，即 $p^*=q$，从而
$h\le\|f-p^*\|_\infty=\|f-q\|_\infty$ 与 $h$ 的定义（下界）不矛盾——
形式化地：把 $q$ 换成 $p^*$ 后同样的论证给出 $\min_j|f(x_j)-p^*(x_j)|\ge h$ 是不可能的，
因为那会超过 $\|f-p^*\|_\infty$；故 $h\le\|f-p^*\|_\infty$。
\end{proof}

这条推论是本章\textbf{最有工程价值}的一句话：它把"最优"变成一个可以当场检验的不等式。
只要交错点上的最小振幅 $h$ 追平了实测最大误差 $\max|e|$，当前多项式就是最优的
（相差因子 $1-\varepsilon$）。Remez 交换算法正是用"换点让 $h$ 尽量大"在跑。

\textbf{本实现的两个可观测结论}（\texttt{test\_remez\_satisfies\_the\_alternation\_theorem}、
\texttt{test\_de\_la\_vallee\_poussin\_certificate}）：交错点个数恒为 $n+2$，
振幅比达到 $1.000000000$；pole 在 $n=8$ 时的最优误差是 $9.8088\times10^{-2}$，
同阶 Chebyshev 插值是 $1.7084\times10^{-1}$，同阶 $L^2$ 投影是 $1.4958\times10^{-1}$
——"最小二乘"在最大范数下输给"最小最大" $1.52$ 倍。

\begin{theorem}[Chebyshev 首一多项式]
\label{thm:monic}
在首一 $n$ 次多项式中，$\|q\|_\infty$ 的最小值是 $2^{1-n}$，且唯一取到于
$q(x)=2^{1-n}T_n(x)$，其零点为 $x_k=\cos\frac{(2k-1)\pi}{2n}$。
\end{theorem}

\begin{proof}
$T_n(\cos\theta)=\cos n\theta$。在 $\theta_j=j\pi/n$ 处 $|T_n|=1$ 且符号交替，
共 $n+1$ 个点，故 $\|2^{1-n}T_n\|_\infty=2^{1-n}$，且它满足交错条件（$n+1$ 个点对
$f\equiv0$ 的 $n-1$ 次逼近……更直接地：$2^{1-n}T_n$ 首一，范数为 $2^{1-n}$）。
唯一性/最优性：设首一 $q$ 满足 $\|q\|_\infty<2^{1-n}$，则 $r=2^{1-n}T_n-q$ 的次数
$\le n-1$（首项相消），而 $r$ 在 $n+1$ 个点 $\theta_j$ 上取值
$2^{1-n}(-1)^j-q(\cos\theta_j)$，符号由 $(-1)^j$ 决定（因为 $|q|<2^{1-n}$），
于是 $r$ 有 $n$ 个零点，只能是 $r\equiv0$，矛盾。
\end{proof}

\textbf{实测}：Chebyshev 零点上的首一多项式范数与 $2^{1-n}$ 的相对偏差 $\le10^{-12}$；
换成等距节点，范数分别是它的 $1.6$、$3.6$、$11.9$、$44.1$ 倍（$n=4,8,12,16$）。
这就是"插值节点不能等距"的定量版本。

\section{误差衰减得有多快}

\begin{theorem}[代数收敛率]
\label{thm:algebraic}
设 $f\in C^\nu[-1,1]$、$f^{(\nu)}$ 有界变差，且 $f$ 在非光滑点处只有有限个折点/跳点。
则 Chebyshev 系数满足 $|c_k|=O(k^{-\nu-1})$，于是 $L^2$ 投影误差
\[
  \|f-P_nf\|_{L^2}=O(n^{-\nu-1/2}).
\]
特别地：$f=|x-a|$（$\nu=1$，折点）给 $O(n^{-3/2})$；$f=\mathrm{sign}(x-a)$
（$\nu=0$，跳点）给 $O(n^{-1/2})$。
\end{theorem}

\begin{proof}[证明要点]
$c_k=\frac{2}{\pi}\int_{-1}^1\frac{f(x)T_k(x)}{\sqrt{1-x^2}}dx=\frac{2}{\pi}\int_0^\pi f(\cos\theta)\cos k\theta\,d\theta$。
对 $\theta$ 做 $\nu$ 次分部积分，每次把 $\cos k\theta$ 变成 $\frac{1}{k}\sin k\theta$，
边界项由 $f$ 的光滑性消去，最后剩下 $k^{-\nu}\int f^{(\nu)}\cos k\theta\,d\theta$；
$f^{(\nu)}$ 的有界变差性（等价于它是某个有限测度的积分）使得最后的积分为 $O(k^{-1})$。
故 $c_k=O(k^{-\nu-1})$。再由 Bessel 恒等式，截断误差
$\|f-P_nf\|_{L^2}^2=\sum_{k>n}c_k^2\cdot\frac{\pi}{1+\delta_{k0}}=O(n^{-2\nu-1})$，
开方即得 $O(n^{-\nu-1/2})$。
\end{proof}

\textbf{实测}（\texttt{test\_algebraic\_rates\_match\_smoothness}）：折点 $|x|$
在 $n\ge12$ 上拟合出的阶是 $1.5106$（理论 $1.5$，$R^2=0.996$）；
跳点 $\mathrm{sign}(x)$ 是 $0.4924$（理论 $0.5$，$R^2=0.993$）。
注意折点那条必须\textbf{在折点处切段求积}：不切段时（固定 40 点求积）量到的阶是 $2.129$，
跳点量到 $0.702$——你测到的是求积误差的阶，不是逼近的阶。

\begin{theorem}[几何收敛与 Bernstein 椭圆]
\label{thm:geometric}
设 $f$ 在包含 $[-1,1]$ 的 Bernstein 椭圆
\[
  E_\rho=\{z:\ |z-1|+|z+1|\le\rho+\rho^{-1}\}\qquad(\rho>1)
\]
内解析，则 $|c_k|\le 2M(\rho)\rho^{-k}$，从而 $\|f-P_nf\|_{L^2}=O(\rho^{-n})$，
半对数图上的斜率是 $-\ln\rho$。若 $f$ 只有极点 $y=\mathrm{i}/\sqrt{s}$（即
$f=1/(1+sx^2)$），则 $\rho=|y|+\sqrt{1+y^2}$。
\end{theorem}

\begin{proof}[证明要点]
把 $x=\frac12(z+z^{-1})$（Joukowski 映射）代入 Chebyshev 系数积分，
$\frac{dz}{z}$ 形式的积分在环域 $\rho^{-1}<|z|<\rho$ 上可以用 Cauchy 积分公式估值：
被积函数 $\frac{f((z+z^{-1})/2)}{z^{k+1}}$ 的模被 $M(\rho)\rho^{-k}$ 控制。
极点在 $z=\mathrm{i}\sqrt s\pm\sqrt{1-s}$ 的映射下给出 $\rho=|y|+\sqrt{1+y^2}$，
其中 $y$ 是极点虚部（$|y|=1/\sqrt s$）。
\end{proof}

\textbf{实测}（\texttt{test\_analytic\_rate\_equals\_minus\_log\_rho}）：
$s=25$ 时 $\rho=1.219803902718557$，理论斜率 $-0.198690$，拟合斜率 $-0.198167$
（相对差 $0.26\%$，$R^2=0.9979$）；$s=4,100$ 的相对差是 $0.26\%,0.60\%$；
但 $s=200$ 时 $\rho=1.0732$、偏差 $2.4\%$——$n\le40$ 还没进渐近区。
"解析函数在半对数图上是直线"这句话，只有在渐近区里成立。

整函数（$f=e^x$）没有有限的 $\rho$：实测 $n\ge15$ 时
$\max_n 2^n\|e_n\|\le7.7\times10^{-10}$，误差比上一阶小得多且比值持续下降——
\textbf{快过任何几何率}。它在半对数图上不是直线，因此"拟合斜率"这件事本身就失效；
再往后误差掉到 $10^{-16}$ 的求值舍入地板，误差曲线变成一条平线，
那时的"交错点"全是噪声（本实现的 Remez 在 $n\ge12$ 上因此报"已到精度极限"）。

\section{四个可视化该怎么看}

\begin{enumerate}
  \item \textbf{projection}：主轴把 $f$ 与"当前 $k$ 维投影"和"目标 $n$ 维投影"画在一起，
        副轴是系数向量与误差曲线。把 \texttt{basis} 从 Legendre 切到 Chebyshev，
        曲线几乎不动（同一个子空间），但副轴的系数向量与 \texttt{meta.gram\_cond} 完全变样；
        切到 monomial，$\kappa$ 从 $1$ 涨到 $10^{17}$。注意误差曲线在端点附近的尖峰——
        它正是"$L^2$ 最好"不等于"处处都好"的证据。
  \item \textbf{convergence\_rates}：双对数轴上三条误差曲线 + 每条自己的实测拟合线与理论参考线。
        理论线只锚定"斜率"，常数项不可知。把 $s$ 从 $25$ 拖到 $200$，理论线变平，
        但实测曲线要更大的 $n$ 才跟上——这就是渐近区的意义。
  \item \textbf{equioscillation}：时间轴是 Remez 的迭代轮次，看误差曲线怎么一步步"压平"、
        交错点怎么滑到极值点上。虚线是同阶 $L^2$ 投影的误差：它在端点更大、
        在中间更小——两条曲线的交叉就是"两种最优"的差别。
  \item \textbf{zeros\_evolution}：横轴是零点位置、纵轴是阶数（分岔图风格），
        叠加理论密度曲线与同阶 Chebyshev 零点的对照。下面板是每个根的 Newton 修正量
        $|p_n/p_n'|$：Legendre 全都贴在 $1.1\times10^{-16}$（地板），
        Laguerre 在 $n=40$ 时最大 $1.2\times10^{-14}$。
  \item 四个可视化都给了 $u/2$ 作为对数轴的地板，并把"被截断的点数"写进
        \texttt{meta.clipped\_*}：$0$ 在对数轴上画不出来，把它当"完美"是最常见的误读。
\end{enumerate}

% 本文件由 tools/build_book.py 从 dist/reports/assertions.json 生成，请勿手改。
\begin{longtable}{@{}p{0.06\textwidth}p{0.32\textwidth}p{0.22\textwidth}p{0.16\textwidth}@{}}
\caption{章节 2.3 数值断言（生成时间 2026-10-19T16:08:50-0400）}\\
\toprule
类别 & 断言 & 理论值 & 实测值 \\
\midrule
\endfirsthead
\toprule
类别 & 断言 & 理论值 & 实测值 \\
\midrule
\endhead
\bottomrule
\endlastfoot

A & max|$\langle$残差, $\varphi$\_\allowbreak{}k$\rangle$|(三函数 $\times$ n=4,8,12) & 0(\allowbreak{}正交) & 5.51e-16 \checkmark \\
A & 逐项明细 & - & pole(\allowbreak{}n=\allowbreak{}4):3.3e-17，pole(\allowbreak{}n=\allowbreak{}8):3.1e-17，pole(\allowbreak{}n=\allowbreak{}12):9.3e-17，cusp(\allowbreak{}n=\allowbreak{}4):2.0e-16，cusp(\allowbreak{}n=\allowbreak{}8):1.6e-16，cusp(\allowbreak{}n=\allowbreak{}12):1.5e-16，entire(\allowbreak{}n=\allowbreak{}4):5.5e-16，entire(\allowbreak{}n=\allowbreak{}8):2.0e-16，entire(\allowbreak{}n=\allowbreak{}12):1.5e-16 \checkmark \\
A & $\|$P$^2$$-$P$\|$$\infty$ 与 5 次多项式的投影误差(n=6,7) & $\le$ 1e-13(\allowbreak{}两者) & 幂等缺陷 2.22e-16，多项式还原误差 4.00e-15 \checkmark \\
A & Clenshaw vs 直接求和 / 基变换 vs 直接展开 & $\le$ 1e-13 & 2.22e-15 与 3.33e-16 \checkmark \\
A & $\|$首一多项式$\|$$\infty$ / 2\textasciicircum{}(1$-$n) $-$ 1(Chebyshev 节点) & 0(\allowbreak{}取到理论最小) & 最大相对偏差 1.69e-14 \checkmark \\
A & 等距节点范数 / Chebyshev 节点范数(n=4,8,12,16) & > 1 且随 n 增长 & 1.6，3.6，11.9，44.1 \checkmark \\
A & 交错点个数 = n+2 且振幅比 $\to$ 1 & n+2 个点、\allowbreak{}比值 $\le$ 1+1e-6 & 最大比值 1.000000000；全部满足 =\allowbreak{} True \checkmark \\
A & 逐例明细 & - & pole(\allowbreak{}n=\allowbreak{}4):6点/\allowbreak{}比值1.000000000/\allowbreak{}收敛；pole(\allowbreak{}n=\allowbreak{}8):10点/\allowbreak{}比值1.000000000/\allowbreak{}收敛；pole(\allowbreak{}n=\allowbreak{}12):14点/\allowbreak{}比值1.000000000/\allowbreak{}收敛；pole(\allowbreak{}n=\allowbreak{}16):18点/\allowbreak{}比值1.000000000/\allowbreak{}收敛；cusp(\allowbreak{}n=\allowbreak{}4):6点/\allowbreak{}比值1.000000000/\allowbreak{}收敛；cusp(\allowbreak{}n=\allowbreak{}8):10点/\allowbreak{}比值1.000000000/\allowbreak{}收敛；cusp(\allowbreak{}n=\allowbreak{}12):14点/\allowbreak{}比值1.000000000/\allowbreak{}收敛；cusp(\allowbreak{}n=\allowbreak{}16):18点/\allowbreak{}比值1.000000000/\allowbreak{}收敛；entire(\allowbreak{}n=\allowbreak{}8):10点/\allowbreak{}比值1.000000080 \checkmark \\
A & max|e| $-$ min|e|(交错集)(Remez 结果) & $\le$ 1e-6$\cdot$max|e|(\allowbreak{}证书成立) & 最大间隙 3.37e-12 \checkmark \\
A & 同一函数：Chebyshev 插值的 L$\infty$ 误差 / Remez 的 L$\infty$ 误差 & > 1(\allowbreak{}非最优多项式的下界与上界对不上) & 1.708356e-01 /\allowbreak{} 9.808814e-02 =\allowbreak{} 1.7417 \checkmark \\
A & 逐例明细 & - & pole(\allowbreak{}n=\allowbreak{}4):h=\allowbreak{}2.171584e-01,\allowbreak{}max=\allowbreak{}2.171584e-01；pole(\allowbreak{}n=\allowbreak{}8):h=\allowbreak{}9.808814e-02,\allowbreak{}max=\allowbreak{}9.808814e-02；pole(\allowbreak{}n=\allowbreak{}12):h=\allowbreak{}4.430538e-02,\allowbreak{}max=\allowbreak{}4.430538e-02；cusp(\allowbreak{}n=\allowbreak{}4):h=\allowbreak{}6.762090e-02,\allowbreak{}max=\allowbreak{}6.762090e-02；cusp(\allowbreak{}n=\allowbreak{}8):h=\allowbreak{}3.468973e-02,\allowbreak{}max=\allowbreak{}3.468973e-02；cusp(\allowbreak{}n=\allowbreak{}12):h=\allowbreak{}2.324733e-02,\allowbreak{}max=\allowbreak{}2.324733e-02 \checkmark \\
A & 尖点 |x| 的 L$^2$ 收敛阶(n$\ge$12) & 1.5 & 1.51106 \checkmark \\
A & 跳点 sign(x) 的 L$^2$ 收敛阶 & 0.5 & 0.490501 \checkmark \\
A & 两条曲线的拟合常数 C 与 R$^2$ & err $\approx$ C$\cdot$n\textasciicircum{}(\allowbreak{}$-$$\alpha$) & cusp C=\allowbreak{}0.676(\allowbreak{}R$^2$=\allowbreak{}0.992324)，jump C=\allowbreak{}1.098(\allowbreak{}R$^2$=\allowbreak{}0.990905) \checkmark \\
A & 半对数斜率 vs $-$ln$\rho$(s=4, $\rho$=1.618034) & 0.481212 & -0.479979(\allowbreak{}R$^2$=\allowbreak{}0.997921) \checkmark \\
A & 半对数斜率 vs $-$ln$\rho$(s=25, $\rho$=1.219804) & 0.198690 & -0.198167(\allowbreak{}R$^2$=\allowbreak{}0.997920) \checkmark \\
A & 半对数斜率 vs $-$ln$\rho$(s=100, $\rho$=1.104988) & 0.099834 & -0.099234(\allowbreak{}R$^2$=\allowbreak{}0.997819) \checkmark \\
A & s=200($\rho$=1.073)时的同一比较 & 偏差变大(\allowbreak{}n$\le$40 还没进渐近区) & -0.068974 vs -0.070652(\allowbreak{}2.38\%) \checkmark \\
A & 拟合明细 & - & s=\allowbreak{}4:-0.479979$\approx$-0.481212(\allowbreak{}0.26\%)；s=\allowbreak{}25:-0.198167$\approx$-0.198690(\allowbreak{}0.26\%)；s=\allowbreak{}100:-0.099234$\approx$-0.099834(\allowbreak{}0.60\%) \checkmark \\
A & max $\rho$\textasciicircum{}n$\cdot$$\|$e\_\allowbreak{}n$\|$(n$\ge$15) & $\to$ 0(\allowbreak{}对每个 $\rho$>1) & $\rho$=\allowbreak{}1.5:4.68e-12，$\rho$=\allowbreak{}2.0:2.63e-09，$\rho$=\allowbreak{}3.0:1.96e-05 \checkmark \\
A & exp(x) 的相邻两点之比 e\_\allowbreak{}\{n+1\}/e\_\allowbreak{}n(只取 e>1e-12 的窗口) & 随 n 严格变小(\allowbreak{}超几何收敛) & n=\allowbreak{}2$\to$3:0.165，n=\allowbreak{}3$\to$4:0.124，n=\allowbreak{}4$\to$5:0.100，n=\allowbreak{}5$\to$6:0.083，n=\allowbreak{}6$\to$7:0.071，n=\allowbreak{}7$\to$8:0.062，n=\allowbreak{}8$\to$9:0.055，n=\allowbreak{}9$\to$10:0.050，n=\allowbreak{}10$\to$11:0.045 \checkmark \\
A & Remez / Chebyshev 插值 与 Remez / L$^2$ 投影 & < 1(\allowbreak{}两者) & 最大比 0.9286 与 0.6569 \checkmark \\
A & 逐例(Remez < cheb < (或) l2) & - & pole(\allowbreak{}n=\allowbreak{}4):2.172e-01 < cheb 4.020e-01 < l2 3.306e-01；pole(\allowbreak{}n=\allowbreak{}8):9.809e-02 < cheb 1.708e-01 < l2 1.496e-01；pole(\allowbreak{}n=\allowbreak{}12):4.431e-02 < cheb 6.922e-02 < l2 6.762e-02；cusp(\allowbreak{}n=\allowbreak{}4):6.762e-02 < cheb 1.232e-01 < l2 1.172e-01；cusp(\allowbreak{}n=\allowbreak{}8):3.469e-02 < cheb 6.697e-02 < l2 6.729e-02；cusp(\allowbreak{}n=\allowbreak{}12):2.325e-02 < cheb 4.613e-02 < l2 4.725e-02；entire(\allowbreak{}n=\allowbreak{}4):5.467e-04 < cheb 6.397e-04 < l2 1.207e-03；entire(\allowbreak{}n=\allowbreak{}8):1.106e-08 < cheb 1.219e-08 < l2 3.136e-08；entire(\allowbreak{}n=\allowbreak{}12):4.041e-14 < cheb 4.352e-14 < l2 1.323e-13 \checkmark \\
B & $\kappa$(G) 随 n 的增长(单项式，n=12$\to$24) & 指数增长(\allowbreak{}实测 1e8 $\to$ 2e17) & 9.96e+07 $\to$ 1.95e+17(\allowbreak{}$\times$2.0e+09)；Legendre n=\allowbreak{}24 的 $\kappa$=\allowbreak{}1.000 \checkmark \\
B & n=24：单项式基投影与良态基投影的函数差 / 正规方程残差 & 函数差 \textasciitilde{} $\kappa$u $\gg$ 残差 & 1.17e-03 vs 1.04e-11 \checkmark \\
B & |投影\_\allowbreak{}Legendre $-$ 投影\_\allowbreak{}Chebyshev| 与 |投影\_\allowbreak{}Legendre $-$ 投影\_\allowbreak{}单项式| & 前者 $\le$1e-13(\allowbreak{}同一子空间)，后者随 $\kappa$ 增大 & 最大 |L$-$C|=\allowbreak{}2.15e-15(\allowbreak{}n$\le$24)；n=\allowbreak{}24 的 |L$-$M|=\allowbreak{}1.17e-03 \checkmark \\
B & 逐阶明细 & - & n=\allowbreak{}12:|L$-$C|=\allowbreak{}4.4e-16,\allowbreak{}|L$-$M|=\allowbreak{}4.7e-11；n=\allowbreak{}20:|L$-$C|=\allowbreak{}5.6e-16,\allowbreak{}|L$-$M|=\allowbreak{}9.1e-06；n=\allowbreak{}24:|L$-$C|=\allowbreak{}2.2e-15,\allowbreak{}|L$-$M|=\allowbreak{}1.2e-03 \checkmark \\
B & 不切段时的实测斜率与理论斜率之差 / 切段后的斜率 & 切段后 |$\alpha$$-$$\nu$$-$1/\allowbreak{}2| $\le$ 0.05；不切段时偏差 $\ge$ 0.15 & 不切段最大偏差 0.629；cusp:不切段 $\alpha$=\allowbreak{}2.129，切段 $\alpha$=\allowbreak{}1.518；jump:不切段 $\alpha$=\allowbreak{}0.702，切段 $\alpha$=\allowbreak{}0.495 \checkmark \\
B & 实测斜率 vs smoothness\_\allowbreak{}table & $\alpha$ 差 $\le$0.1，斜率差 $\le$1\% & cusp:实测 1.511 vs 表 1.5；jump:实测 0.491 vs 表 0.5；pole:实测 -0.197876 vs 表 -0.198690 \checkmark \\
B & 跳点 n=6：converged / 返回的 max|e| / 历史最优 & converged=\allowbreak{}False 且 max|e| =\allowbreak{} min(\allowbreak{}历史) & converged=\allowbreak{}False，max|e|=\allowbreak{}1.000e+00，历史最优=\allowbreak{}1.000e+00 \checkmark \\
B & 同一函数：L$^2$ 投影误差仍在按 n\textasciicircum{}($-$1/2) 下降 & 单调下降 & 5.303e-01，3.867e-01，2.777e-01，2.279e-01 \checkmark \\
B & 帧契约检查(4 个 viz，含时间轴有效性) & requiredAux 齐备、\allowbreak{}数组全有限、\allowbreak{}meta 只放有限标量、\allowbreak{}scrub 帧随时间变化 & projection: 8 数组，meta 12 项；convergence\_rates: 11 数组，meta 18 项；equioscillation: 9 数组，meta 16 项；zeros\_evolution: 9 数组，meta 14 项 \checkmark \\
A & 递推 vs 闭式(n$\le$12 / n$\le$16) & $\le$ 1e-11(\allowbreak{}Legendre)、\allowbreak{}$\le$ 1e-13(\allowbreak{}Chebyshev)、\allowbreak{}$\le$ 1e-8(\allowbreak{}Laguerre,\allowbreak{} x$\le$10) & Legendre 3.01e-13，Chebyshev 6.08e-15，Laguerre 6.24e-10 \checkmark \\
A & Laguerre 递推误差的窗口依赖(x 上限 5/10/20) & 绝对误差随 x 上限增长(\allowbreak{}|L\_n| 自身量级被放大) & 9.56e-12 /\allowbreak{} 6.24e-10 /\allowbreak{} 1.25e-07 \checkmark \\
A & 闭式 vs SciPy(n$\le$12) & $\le$ 1e-12 & Legendre 1.51e-13，Chebyshev 3.55e-15 \checkmark \\
A & max|$\langle$p\_\allowbreak{}i,p\_\allowbreak{}j$\rangle$\_\allowbreak{}w $-$ $\delta$\_\allowbreak{}ij| & $\le$ 1e-13(\allowbreak{}舍入量级) & 7.59e-14(\allowbreak{}legendre(\allowbreak{}n$\le$6)=\allowbreak{}1.8e-15，legendre(\allowbreak{}n$\le$12)=\allowbreak{}2.3e-15，legendre(\allowbreak{}n$\le$20)=\allowbreak{}1.7e-15，chebyshev(\allowbreak{}n$\le$6)=\allowbreak{}7.5e-16，chebyshev(\allowbreak{}n$\le$12)=\allowbreak{}1.7e-15，chebyshev(\allowbreak{}n$\le$20)=\allowbreak{}2.4e-15，laguerre(\allowbreak{}n$\le$6)=\allowbreak{}1.9e-15，laguerre(\allowbreak{}n$\le$12)=\allowbreak{}7.6e-14，laguerre(\allowbreak{}n$\le$20)=\allowbreak{}4.0e-14) \checkmark \\
A & Gram 矩阵非对角元(legendre, n$\le$10) & $\approx$ 0 & 5.72e-16(\allowbreak{}对角元偏离 1 的最大值 4.44e-16) \checkmark \\
A & k$\le$2n$-$1 的最大矩误差 / k=2n 处的误差 & 前者 $\le$ 1e-13，后者 $\ge$ 1e-10 & 精确段 3.16e-14，失效点最小值 7.20e-10 \checkmark \\
A & 逐族明细(精确段/失效点) & 2n$-$1 是硬边界 & legendre(\allowbreak{}n=\allowbreak{}4):4.4e-16/\allowbreak{}1.2e-02，chebyshev(\allowbreak{}n=\allowbreak{}4):1.1e-16/\allowbreak{}2.5e-02，laguerre(\allowbreak{}n=\allowbreak{}4):8.9e-16/\allowbreak{}1.4e-02，legendre(\allowbreak{}n=\allowbreak{}8):8.9e-16/\allowbreak{}4.7e-05，chebyshev(\allowbreak{}n=\allowbreak{}8):2.2e-16/\allowbreak{}9.6e-05，laguerre(\allowbreak{}n=\allowbreak{}8):3.4e-15/\allowbreak{}7.8e-05，legendre(\allowbreak{}n=\allowbreak{}16):4.4e-16/\allowbreak{}7.2e-10，chebyshev(\allowbreak{}n=\allowbreak{}16):6.1e-16/\allowbreak{}1.5e-09，laguerre(\allowbreak{}n=\allowbreak{}16):3.2e-14/\allowbreak{}1.7e-09 \checkmark \\
A & |$\Sigma$w $-$ (b$-$a)|(n$\le$64，三族) & $\le$ 1e-12 & 3.33e-13；最小权重 2.09e-101 > 0 \checkmark \\
A & 逐族明细 & Chebyshev 的权重是常数 $\pi$/\allowbreak{}n(\allowbreak{}精确) & legendre(\allowbreak{}n=\allowbreak{}4):4.4e-16，chebyshev(\allowbreak{}n=\allowbreak{}4):0.0e+00，laguerre(\allowbreak{}n=\allowbreak{}4):8.9e-16，legendre(\allowbreak{}n=\allowbreak{}16):4.4e-16，chebyshev(\allowbreak{}n=\allowbreak{}16):0.0e+00，laguerre(\allowbreak{}n=\allowbreak{}16):3.2e-14，legendre(\allowbreak{}n=\allowbreak{}32):8.9e-16，chebyshev(\allowbreak{}n=\allowbreak{}32):0.0e+00，laguerre(\allowbreak{}n=\allowbreak{}32):2.4e-14，legendre(\allowbreak{}n=\allowbreak{}64):4.4e-16，chebyshev(\allowbreak{}n=\allowbreak{}64):0.0e+00，laguerre(\allowbreak{}n=\allowbreak{}64):3.3e-13 \checkmark \\
A & family\_\allowbreak{}norm vs 闭式 / vs 数值积分 & $\le$ 1e-12 & 闭式差 0.00e+00，数值积分差 4.44e-16 \checkmark \\
A & 前几项范数 & 2/\allowbreak{}(\allowbreak{}2k+1)、\allowbreak{}$\pi$/\allowbreak{}$\pi$2、\allowbreak{}1 & legendre(\allowbreak{}k=\allowbreak{}0):2.000000$\approx$2.000000，legendre(\allowbreak{}k=\allowbreak{}1):0.666667$\approx$0.666667，legendre(\allowbreak{}k=\allowbreak{}2):0.400000$\approx$0.400000，chebyshev(\allowbreak{}k=\allowbreak{}0):3.141593$\approx$3.141593，chebyshev(\allowbreak{}k=\allowbreak{}1):1.570796$\approx$1.570796，chebyshev(\allowbreak{}k=\allowbreak{}2):1.570796$\approx$1.570796，laguerre(\allowbreak{}k=\allowbreak{}0):1.000000$\approx$1.000000，laguerre(\allowbreak{}k=\allowbreak{}1):1.000000$\approx$1.000000，laguerre(\allowbreak{}k=\allowbreak{}2):1.000000$\approx$1.000000 \checkmark \\
A & max|V $-$ 正确方向的零点个数|(探针避开零点) & 0(\allowbreak{}精确相等) & 0(\allowbreak{}共核对 369 个探针点) \checkmark \\
A & 把方向搞反时的最大偏差 & > 0(\allowbreak{}方向必须按族的归一化端点选) & 12 \checkmark \\
A & Gram--Schmidt vs 解析 Legendre(Gauss--Legendre 测度，200 点) & $\le$ 1e-12 & 5.46e-14；GS 自报缺陷 4.44e-16 \checkmark \\
A & 换成反正弦测度后 GS 得到的其实是 Chebyshev 族 & $\le$ 1e-12(\allowbreak{}同一个算法，测度变了族就变了) & 1.89e-15 \checkmark \\
A & 零点处 |T\_\allowbreak{}n| / 极点处 ||T\_\allowbreak{}n|$-$1| & 0 与 0(\allowbreak{}闭式) & 1.23e-14 与 0.00e+00 \checkmark \\
A & 极点处符号严格交替(n=4,8,16,32) & 符号逐个交替 & [True,\allowbreak{} True,\allowbreak{} True,\allowbreak{} True] \checkmark \\
B & 根数/升序/收敛(三族，n$\le$32) & 根数 =\allowbreak{} n，间距 > 0，全部收敛 & 全部满足 =\allowbreak{} True；最慢单根迭代次数 4 \checkmark \\
B & 各族的逐根最大迭代次数 & $\le$ 60 & legendre:maxIter=\allowbreak{}4，chebyshev:maxIter=\allowbreak{}1，laguerre:maxIter=\allowbreak{}2 \checkmark \\
B & Laguerre：残差 / 根误差(n$\le$40) & 残差 $\gg$ 根误差(\allowbreak{}两者相差 $\ge$1e25) & 残差 5.63e+14，根误差 1.20e-14，比值 4.70e+28 \checkmark \\
B & 逐阶明细 & n 越大残差越大 & n=\allowbreak{}8:res=\allowbreak{}7.3e-12,\allowbreak{}err=\allowbreak{}1.1e-15，n=\allowbreak{}16:res=\allowbreak{}1.9e-05,\allowbreak{}err=\allowbreak{}3.0e-15，n=\allowbreak{}32:res=\allowbreak{}3.4e+07,\allowbreak{}err=\allowbreak{}6.8e-15，n=\allowbreak{}40:res=\allowbreak{}5.6e+14,\allowbreak{}err=\allowbreak{}1.2e-14 \checkmark \\
B & 每两个相邻低阶零点之间的高阶零点个数 $-$ 1 & 0(\allowbreak{}严格交错) & 最大偏离 0(\allowbreak{}共核对 132 个区间) \checkmark \\
B & 零点在正交区间内 + 零点处的 |p\_\allowbreak{}n|(递推回代) & 区间内 + |p\_n| $\le$ 1e-11 & 区间检查 =\allowbreak{} True；Legendre/\allowbreak{}Chebyshev 残差 1.35e-14 \checkmark \\
B & 同一个零点用闭式展开式回代得到的残差 & $\gg$ 1e-11(\allowbreak{}闭式自身的大数相消，不是求根不准) & 2.28e-06(\allowbreak{}n=\allowbreak{}32 的 Legendre) \checkmark \\
B & 权重与闭式公式的相对差(Legendre \& Laguerre，n$\le$32) & $\le$ 1e-12 & 2.02e-16 \checkmark \\
B & Legendre 权重逐阶明细 & - & n=\allowbreak{}4:1.7e-16，n=\allowbreak{}8:1.5e-16，n=\allowbreak{}16:1.5e-16，n=\allowbreak{}32:1.4e-16 \checkmark \\
B & quadrature() 与 $\Sigma$w$\cdot$f 的相对差 / 权函数全正 & 0 与 True & 0.00e+00 与 True \checkmark \\
B & Laguerre 权函数积分 $\int$e\textasciicircum{}\{$-$x\}dx & 1 & 0.9999999999999850 \checkmark \\
B & 精确矩 vs 求积(k=6，三族) & $\le$ 1.00000e-13 & $5.684\times 10^{-15}$ \checkmark \\
B & G 的对称性/对角/非对角偏离 & $\le$ 1e-13 & 非对称 1.1e-16，对角偏离 1.8e-15，非对角 1.4e-15 \checkmark \\
\end{longtable}


\section{它会怎样骗你（本章版）}

\begin{enumerate}
  \item \textbf{把 $L^2$ 最好当成处处都好}：$1/(1+25x^2)$ 在 $n=8$ 时 $L^2$ 投影的最大误差是
        $1.50\times10^{-1}$，而 Remez 是 $9.81\times10^{-2}$——差 $1.52$ 倍，
        而且在端点附近差得更多。
  \item \textbf{以为"换个基结果就不同"}：三个基底张成同一个子空间，投影作为函数相同
        （Legendre 与 Chebyshev 的差 $\le2.2\times10^{-15}$）。变的是系数与条件数。
  \item \textbf{把残差小当成解对}：单项式基 $n=24$ 时正规方程残差 $1.0\times10^{-11}$，
        看起来完美，而 $\kappa u=21.7$，投影函数已经错了 $1.2\times10^{-3}$。
  \item \textbf{用残差当求根误差}：Laguerre $n=40$ 的 $|L_n(x_k)|$ 是 $5.6\times10^{14}$，
        写成误差就是"算法崩了"；Newton 修正量 $|L_n/L_n'|$ 只有 $1.2\times10^{-14}$。
        量误差前先问一句：这个量跟多项式的尺度有关吗？
  \item \textbf{用闭式展开式在高阶回代}：$P_{32}$ 的显式展开式自身只剩 $10^{-6}$ 精度
        （大数相消），拿它当"真值"会以为求根错了。递推才是稳定的那条路。
  \item \textbf{不切段就拟合收敛率}：不切段时 $|x|$ 的阶是 $2.13$、$\mathrm{sign}(x)$ 是 $0.70$，
        切段后才是 $1.52$、$0.495$。误差曲线的斜率随时可以是求积的斜率。
  \item \textbf{用"交错点不再移动"当停机判据}：浮点噪声下点集会在两三个状态之间循环，
        其中既有最优解、也有误差大一倍的解。必须用 de la Vall\'ee Poussin 下界
        追平实测上界来判断（本实现若拿不到证书就返回历史最优迭代并如实报
        \texttt{converged=False}）。
  \item \textbf{在舍入地板上数交错点}：$e^x$ 在 $n\ge12$ 时误差 $\sim10^{-14}$，
        "数出 $n+2$ 个等幅点"毫无意义；本实现把这种情形记为"到精度极限"。
  \item \textbf{把 $\sum w_k=b-a$ 当成"公式对了"}：权重和对了但权重公式错一个因子，
        在低阶上可能完全看不出来；所以还必须在 $k=2n$ 处看到精度\textbf{当场失效}。
  \item \textbf{把时间轴当时间}：它是"已展开的基函数个数"（projection）、
        Remez 迭代轮次（equioscillation）、最高阶数（zeros\_evolution），
        拖到头就夹住，真实值看 \texttt{meta.step}。
  \item \textbf{把理论参考线当成拟合结果}：图上的理论线只固定了\textbf{斜率}，
        常数项是锚在第一个实测点上的；两条线重合说明斜率对，不说明常数对。
  \item \textbf{把零点画成等距的}：Legendre $n=24$ 时相邻根间距最大比最小是 $6.26$ 倍，
        Laguerre $n=40$ 时是 $92.3$ 倍；它们服从反正弦律 / Marchenko--Pastur 律。
\end{enumerate}
